\begin{problem}[Reconstruct the sheaf from principal-open data]\label{prob:vi17-p18}
Suppose a sheaf of rings $\mathcal F$ on $\Spec A$ is equipped with compatible isomorphisms $\mathcal F(D(f))\cong A_f$ for every $f\in A$.  Prove $\mathcal F\cong\OO_{\Spec A}$.
\end{problem}

\ifdefined\FullProblemDossiers
\paragraph{Why this problem matters.}
Principal opens form a basis, so the structure sheaf is determined by its localization data there.

\paragraph{Useful ideas.}
A sheaf is determined by its restriction to a basis if the restriction maps and overlap identifications agree.

\paragraph{Guided hints.}
\begin{enumerate}
\item Every open set is covered by principal opens.
\item Map sections locally using the given isomorphisms.
\item Compatibility on basis overlaps lets the local images glue.
\end{enumerate}

\begin{solution}
For an open $U$, cover $U$ by principal opens $D(f_i)$.  A section of $\mathcal F(U)$ restricts to compatible sections in $\mathcal F(D(f_i))\cong A_{f_i}=\OO_X(D(f_i))$.  These compatible structure-sheaf sections glue uniquely to a section of $\OO_X(U)$.  The construction is independent of the chosen cover because any two covers admit a principal refinement.  Reversing the given basis isomorphisms constructs the inverse map.  Compatibility with restrictions follows from the same refinement argument.
\end{solution}

\paragraph{What happened mathematically.}
The entire sheaf is encoded by its values on the distinguished basis plus compatible restriction maps.

\paragraph{Extensions.}
Use this principle in VI/18 to identify $\OO_X|_{D(f)}$ with the structure sheaf of $A_f$.

\paragraph{Related problems.}
VI.17.P17 and VI.17.P19 develop neighboring aspects of the same localization-sheaf calculus.

\paragraph{Provenance.}
\textbf{New synthesis.  It prepares but does not consume AG3 Problem 11.}
\else
\paragraph{Provenance.} \textbf{New synthesis.  It prepares but does not consume AG3 Problem 11.}
\fi
